\uuid{mjoB}
\exo7id{1799}
\titre{exo7 1799}
\auteur{ridde}
\organisation{exo7}
\datecreate{1999-11-01}
\isIndication{false}
\isCorrection{false}
\chapitre{Fonction de plusieurs variables}
\sousChapitre{Différentiabilité}
\cpge{mpsi-mp2i.fonctions-deux-variables.b, pcsi.fonctions-deux-variables.b, ptsi.fonctions-deux-variables.b, mpsi-mp2i.limites-continuite.b, pcsi.limites-continuite.b, ptsi.limites-continuite.b, bcpst1.fonctions-deux-variables.b, bcpst1.fonctions-deux-variables.c, mp-mpi.calcul-differentiel-optimisation.a, pt.fonctions-deux-trois-variables.g}
\module{Analyse}
\niveau{L2}
\difficulte{}

\contenu{
\question{
Soit $f : \begin{cases} \Rr^{2} \rightarrow \Rr \\ (x, y) \mapsto \frac{\sin (xy)}
{\left|x\right| + \left|y\right|} \text{ si } (x, y) \neq (0,0)\\
 (0, 0) \mapsto 0\\ \end{cases}$.\\
\'Etudier la continuit\'e de $f$ et l'existence des d\'eriv\'ees partielles.$f$ est-elle
 $C^{1}$ ?
}
}
